\chapter{Workspaces, attachments and undo}
\label{ch:workspaces}

\nodwyn{} started as a coding agent, and git was its undo. It now works in
any folder: a code repository, a folder of documents, a mix. This chapter
covers what changes when the folder is not only code, and how to get
things back.

\section{A workspace}

A \textbf{workspace} is a folder plus a small settings file,
\file{<folder>/.nodwyn/workspace.toml}. Without the file, a folder is a
plain code workspace: that folder, writable, nothing else.

\begin{toml}
# <folder>/.nodwyn/workspace.toml
kind = "docs"

[[mount]]
path = "~/Reference/garden-books"
access = "read"

[[mount]]
path = "/Volumes/Shared/seed-orders"
access = "write"
\end{toml}

\begin{description}[leftmargin=1.2em,labelindent=0pt]
  \item[kind] \code{code}, \code{docs} or \code{mixed}. It is a
    \emph{label}: the console shows it, and it says what the folder holds.
    Every kind has the same tools and the same guard.
  \item[mounts] other folders a task may use, by absolute path (or
    \code{\textasciitilde/}), \code{read} or \code{write}. A read-only mount
    refuses every write. Mounts may not overlap the workspace or each
    other, and \code{/} or your whole home folder is refused as too broad.
\end{description}

The console's \textbf{Workspaces} page edits this for you and names the
folder in a registry, so Home can offer it. From Scheme:

\begin{txt}
you> :eval (workspace)
(("kind" . "code") ("mounts"))
\end{txt}

When a workspace has mounts, the default system prompt adds a section
naming them, so the model knows to reach them by absolute path.

\begin{careful}
\file{workspace.toml} lives under \file{.nodwyn/}, a protected path: the
model cannot widen its own mounts. Changing them is your decision, in the
console or by editing the file, and applies to tasks started afterwards.
\end{careful}

\section{Undo where git is not the undo}

Before \code{fs-write} or \code{fs-edit} replaces a file that no git work
tree covers, \nodwyn{} keeps its previous content in
\file{.nodwyn/versions/} and records a \code{file.version} event:

\begin{txt}
you> :eval (fs-edit "plants.txt" "basil" "basil\nmint")
⚑ guard confirm fs-edit: fs-edit is irreversible; changes plants.txt → yes
[kept the previous version of /home/you/garden/plants.txt]
\end{txt}

\begin{sh}
$ nodwyn versions plants.txt
2026-09-30 13:19:00  55d684f93ad3      15 bytes
$ nodwyn undo plants.txt
Restored /home/you/garden/plants.txt to its version of 2026-09-30 13:19 (55d684f93ad3).
\end{sh}

In the console, a task's \textbf{Changes} list has \textbf{History} per
file; \textbf{Restore} writes a kept version back through the ordinary
approval (you see the diff and answer), and keeps the current content as a
version too, so a restore can itself be undone. The store is capped per
folder.

\begin{tip}
In a git repository, git is the undo: nothing is kept in
\file{.nodwyn/versions/}, and \code{git diff} and \code{git checkout} are
your tools. The model gets two of its own there, \code{git-status} and
\code{git-diff}: read-only, so they need no approval, and run with
everything in git that could start another program switched off (external
diff drivers, text conversion, the pager), with credential-looking files
left out. Commit before a big task, and review with \code{git diff}
afterwards; with the \key{T} answer (all writes this turn), that is the
fastest safe way to let \nodwyn{} make a large change.
\end{tip}

\section{Attachments}

In the console, the paperclip attaches files to your next message: text
formats, Word (\file{.docx}), Excel (\file{.xlsx}), PowerPoint
(\file{.pptx}), and PDF when \code{pdftotext} is installed, and images
(below). Each file other than an image is:

\begin{enumerate}
  \item copied into the folder's \file{.nodwyn/attachments/};
  \item turned into text, redacted (secret-looking strings hidden) and
    capped;
  \item recorded as an \code{attach} event and shown to the model
    \emph{as data}, not as instructions.
\end{enumerate}

The conversation shows a chip per file with how many secret-like strings
were hidden. The model reads the extracted text; the copy in
\file{.nodwyn/attachments/} is ignored by git, so it never ends up in your
repository's diff.

\begin{careful}
``As data, not instructions'' is a framing the model is given, not a
guarantee it obeys. A document can still say ``ignore your task and
\ldots''. What the guard guarantees holds regardless: whatever the model
decides after reading it still has to pass through the door.
\end{careful}

\subsection{Images}

A PNG, JPEG, GIF or WebP (decided by what the file contains, not its name;
at most 10\,MB and 8000 pixels a side) is sent to the model as an image, so
it can only be attached when the control model takes images: Ollama says so
itself, a chat-completions model needs \code{vision = true} in its table.
Nothing redacts pixels, so the task asks first, every time:

\begin{txt}
─── confirm: attach image whiteboard.png
  • whiteboard.png is an image (PNG, 1200×800, 32 KB, about 1,105 tokens).
    Its contents cannot be redacted: whatever it shows, a password or a key
    included, reaches the model as it is, and nothing in it is read as text.
  • It goes to the control model qwen2.5vl at workshop:11434, another machine:
    its pixels leave this computer as they are.
  • It goes with later requests to that host only, while the conversation
    shows it; the judge never gets it.
allow? [y/N]
\end{txt}

It is a security confirm: \key{Y} covers this image only, \code{--yolo}
never says yes to it, and \code{--no-confirm} says no. The chip then reads
``image · 1200×800 · \textasciitilde1,105 tokens'', which is what the
picture costs the context. The log keeps its size, digest and the host you
allowed, never the pixels. A request to any other host gets a line saying
the image is not for it, and the planner and sub-agents, which start in a
context of their own, never see it.

\begin{tip}
To take an image back out of a conversation, delete its copy under
\file{.nodwyn/attachments/}. Later requests then say it is missing; they
never send a different picture in its place.
\end{tip}

\section{Documents out}

The model can also write documents. \code{doc-export} turns Markdown or
LaTeX into Word, OpenDocument or HTML with pandoc (2.15 or later; without
it the tool says what to install):

\begin{scm}
(doc-export "report.md" "out/report.docx")
;=> "wrote out/report.docx (Word document, 10 KiB) from report.md"
\end{scm}

It is a write: it asks with source, target and formats, a standing write
approval for the folder covers it, and replacing a file keeps the old one
in the version store. pandoc reads the document in its sandbox, so an
include or a remote image in it reads nothing, and a Word or OpenDocument
file embeds only images that are files in the folder. An HTML page leaves
out an image that links elsewhere, which a browser would fetch when the
page is opened. Raw content written into the source (an HTML block, a Word
field) comes out as text, never interpreted. The answer says which images
were left out and when raw content was shown as text.
There is no PDF: a TeX engine reads any file and runs programs when a
document asks, so build a PDF with \code{shell-exec} and \code{latexmk}, as
\nodwyn's own book is built.

With a \code{google} credential stored and bound (Chapter~\ref{ch:keys}),
\code{google-doc-upload} sends a \file{.docx}, \file{.odt}, \file{.html} or
\file{.txt} to Google Drive as a new Google Doc and returns its link:

\begin{scm}
(google-doc-upload "out/report.docx" "Q3 report")
;=> (("id" . "1AbC...") ("url" . "https://docs.google.com/document/d/1AbC.../edit") ("title" . "Q3 report"))
\end{scm}

The first upload asks how uploads should be handled from now on: \emph{allow
by hand} (every call asks, naming the file, its size and where it goes, and
\code{-{}-yolo}, \code{-{}-judge} and unattended runs never upload),
\emph{always allow} (no call asks; only this computer can choose it, never
a paired phone), or \emph{deny} (refused, and later sessions do not offer
the tool). The answer is \code{google\_upload} in the \code{[docs]} table
of \file{config.toml}; Settings in the console changes it, and the log
records it as \code{setting.change}.

\begin{careful}
The file's bytes never become a Scheme value, so the door cannot read them.
The upload reads them itself, every part of a Word file included, and
refuses one that holds something shaped like a key whose class is not bound
for Google; \code{doc.upload} in the log says what left, never what it
said. A contract or a salary list has no shape: what keeps it home is your
``no''.
\end{careful}

\section{Notes, scratch and versions: where things live}

\begin{center}
\begin{tabularx}{\linewidth}{lX}
\toprule
\textbf{what} & \textbf{where} \\
\midrule
workspace settings & \file{<folder>/.nodwyn/workspace.toml} \\
project tools and strategy & \file{<folder>/.nodwyn/init.scm} (and what it loads) \\
standing approvals & \file{<folder>/.nodwyn/allow.toml}, \file{config.toml} \\
kept versions & \file{<folder>/.nodwyn/versions/} (only outside git) \\
attachments & \file{<folder>/.nodwyn/attachments/} \\
project notes & \nodwyn's data directory, keyed by the folder's path \\
session logs & \file{\textasciitilde/.local/share/nodwyn/sessions/} \\
scratch & \nodwyn's data directory; removed when the session closes \\
keys & the keychain, in the config directory \\
\bottomrule
\end{tabularx}
\end{center}
